% Libro 2 · Mundo 4 · Conteo y combinatoria (app: pu-conteo)
\mundo{Conteo y combinatoria}{pu-conteo}{Contar parece fácil hasta que las posibilidades son miles. La combinatoria es el
arte de contar sin hacer la lista: multiplicar decisiones, decidir si el orden importa, separar por casos o contar lo
contrario. Cinco ideas que resuelven casi todo.}{%
\item \textbf{4.1} Principio multiplicativo y árboles
\item \textbf{4.2} Permutaciones: cuando el orden importa
\item \textbf{4.3} Combinaciones: cuando el orden no importa
\item \textbf{4.4} Conteo por casos y por complemento
\item \textbf{4.5} Contar figuras y caminos}

% ------------------------------------------------------------------
\leccion{pu-co-1}{Principio multiplicativo y árboles}{Juego del árbol que se ramifica.}

\subsection{La idea: decisiones en cadena}
Si una tarea se hace en pasos, y el primer paso tiene $m$ opciones y el segundo $n$ (sin importar qué se eligió en el
primero), la tarea completa tiene $m \cdot n$ maneras. Es el \clave{principio multiplicativo}, y un \clave{árbol} lo
muestra:

\figura{arbolconteo}{Cada camisa se combina con cada pantalón: las ramas se multiplican.}

\subsubsection{Sumar o multiplicar}
Se \emph{multiplica} cuando hay que tomar una decisión \emph{y} luego otra. Se \emph{suma} cuando se elige entre caminos
excluyentes: \emph{o} esto \emph{o} aquello. Para ir de Medellín a Bogotá hay 3 vuelos o 5 buses: $3 + 5 = 8$ maneras.

\subsubsection{Empezar por la casilla con restricción}
Cuando un paso tiene restricciones (la primera cifra no puede ser 0, el número debe ser par), se llena primero la casilla
más restringida.

\begin{regla}
\begin{itemize}
  \item «Y luego»: se multiplica. «O» excluyente: se suma.
  \item Llena primero las casillas con restricciones.
  \item Con repetición permitida, cada casilla conserva todas sus opciones: $n^k$.
\end{itemize}
\end{regla}

\subsection{Más ejemplos, en corto}
\begin{ejemplo}[Placas]
Una placa tiene 3 letras (de 26) y 3 dígitos. ¿Cuántas placas distintas hay?
\solucion $26^3 \cdot 10^3 = 17\,576\,000$.
\end{ejemplo}

\begin{ejemplo}[Pares de tres cifras distintas]
¿Cuántos números pares de tres cifras distintas hay?
\solucion Restricciones: la última par, la primera no cero. Caso última $= 0$: primera 9 opciones, del medio 8: 72. Caso
última $\in \{2, 4, 6, 8\}$: 4 opciones; primera: 8 (ni 0 ni la última); del medio: 8. $4 \cdot 8 \cdot 8 = 256$. Total
$328$.
\end{ejemplo}

\begin{ejemplo}[Menú]
Un restaurante ofrece 4 entradas, 5 platos fuertes y 3 postres. Un menú tiene un plato fuerte y, opcionalmente, una entrada
y un postre. ¿Cuántos menús hay?
\solucion Entrada: 4 o ninguna (5 opciones). Fuerte: 5. Postre: 3 o ninguno (4). $5 \cdot 5 \cdot 4 = 100$.
\end{ejemplo}

\begin{trampa}
Olvidar el caso «ninguno» cuando algo es opcional, o empezar por una casilla sin restricción y luego no saber cuántas
opciones quedan para la restringida.
\end{trampa}

\begin{reto}
¿Cuántos números de tres cifras \emph{distintas} se pueden formar con los dígitos 0, 1, 2, 3 y 4?
\opciones{48}{60}{100}{64}
\end{reto}

\begin{retoresuelto}
Primera cifra: 4 opciones (no el 0). Segunda: 4 (ya se usó una, pero ahora el 0 vale). Tercera: 3. $4 \cdot 4 \cdot 3 =
48$. \textbf{Opción A.} La B ($5 \cdot 4 \cdot 3$) cuenta números que empiezan en 0.
\end{retoresuelto}

\taller
\piso{Entrena}{multiplicar}
\begin{ejercicios}
\item ¿Cuántas contraseñas de 4 dígitos hay (se pueden repetir)? \sol{$10^4 = 10\,000$.}
\item ¿Cuántos números de 3 cifras (pueden repetirse) hay que usen solo dígitos impares? \sol{$5^3 = 125$.}
\item Una moneda se lanza 5 veces. ¿Cuántas secuencias de caras y sellos hay? \sol{$2^5 = 32$.}
\end{ejercicios}
\piso{Como en el examen}{opción múltiple}
\begin{ejercicios}
\item ¿Cuántos números de 4 cifras distintas se forman con 1, 2, 3, 4, 5, 6 si deben empezar en 5?
\opciones{60}{120}{20}{360}
\sol{A. $1 \cdot 5 \cdot 4 \cdot 3$.}
\item Entre las ciudades A y B hay 3 caminos y entre B y C hay 4. ¿De cuántas maneras se va de A a C y se regresa sin repetir ningún camino?
\opciones{144}{72}{24}{12}
\sol{B. Ida $3 \cdot 4 = 12$; vuelta $3 \cdot 2 = 6$: 72.}
\end{ejercicios}

\bitacora

% ------------------------------------------------------------------
\leccion{pu-co-2}{Permutaciones}{Juego de los libros que se ordenan en el estante.}

\subsection{La idea: ordenar todos}
Para ordenar $n$ objetos distintos hay $n$ opciones para el primer lugar, $n - 1$ para el segundo, y así hasta 1. Eso es
el \clave{factorial}: $n! = n \cdot (n - 1) \cdots 2 \cdot 1$. Cinco libros se ordenan de $5! = 120$ maneras.

\subsubsection{Ordenar solo algunos}
Si de $n$ objetos se ordenan solo $k$ (un podio de 3 entre 8 corredores), hay $n(n-1)\cdots(n-k+1)$ maneras:
$8 \cdot 7 \cdot 6 = 336$. Se escribe $P(n, k) = \tfrac{n!}{(n-k)!}$.

\subsubsection{Bloques y repetidos}
\begin{itemize}
  \item Si dos deben ir juntos, se pegan en un \clave{bloque}: se ordenan los bloques y luego se ordena dentro del bloque.
  \item Si hay objetos repetidos, se divide por las maneras de ordenar los repetidos: las letras de ANA se ordenan de
  $\tfrac{3!}{2!} = 3$ maneras (ANA, AAN, NAA).
  \item En una mesa redonda, $n$ personas se sientan de $(n - 1)!$ maneras (se fija a una).
\end{itemize}

\begin{regla}
\begin{itemize}
  \item $n!$ maneras de ordenar $n$ objetos distintos. $0! = 1$.
  \item $P(n, k) = n(n - 1)\cdots(n - k + 1)$.
  \item Con repetidos: $\dfrac{n!}{a!\, b!\, \cdots}$. Circular: $(n - 1)!$.
\end{itemize}
\end{regla}

\subsection{Más ejemplos, en corto}
\begin{ejemplo}[Anagramas]
¿Cuántas palabras (con o sin sentido) se forman con las letras de PAPAYA?
\solucion 6 letras: P dos veces, A tres veces, Y una. $\tfrac{6!}{2!\,3!} = \tfrac{720}{12} = 60$.
\end{ejemplo}

\begin{ejemplo}[Separados]
¿De cuántas maneras se sientan 5 personas en fila si Ana y Beto \emph{no} deben quedar juntos?
\solucion Total $120$; juntos $48$ (el reto). Separados: $72$.
\end{ejemplo}

\begin{ejemplo}[Cargos]
En un grupo de 10 se elige presidente, secretario y tesorero. ¿De cuántas maneras?
\solucion Los cargos son distintos: importa el orden. $10 \cdot 9 \cdot 8 = 720$.
\end{ejemplo}

\begin{trampa}
Olvidar ordenar dentro del bloque: Ana-Beto y Beto-Ana son distintos.
\end{trampa}

\begin{reto}
¿De cuántas maneras se pueden sentar 5 personas en fila si Ana y Beto deben quedar juntos?
\opciones{24}{48}{120}{96}
\end{reto}

\begin{retoresuelto}
El bloque \fbox{Ana Beto} y las otras 3 personas son 4 objetos: $4! = 24$ órdenes. Dentro del bloque, 2 órdenes. Total
$48$. \textbf{Opción B.}
\end{retoresuelto}

\taller
\piso{Entrena}{ordenar}
\begin{ejercicios}
\item $\dfrac{8!}{6!}$ \sol{56.}
\item ¿Cuántos anagramas tiene la palabra LIBRO? ¿Y ARROZ? \sol{120 y 60.}
\item ¿De cuántas maneras se sientan 6 personas en una mesa redonda? \sol{$5! = 120$.}
\end{ejercicios}
\piso{Como en el examen}{opción múltiple}
\begin{ejercicios}
\item ¿De cuántas maneras pueden llegar a la meta 6 corredores si se premian el 1.º, el 2.º y el 3.º lugar?
\opciones{20}{120}{720}{216}
\sol{B. $6 \cdot 5 \cdot 4$.}
\item Tres libros de matemáticas y dos de física se ordenan en un estante con los de cada materia juntos. ¿De cuántas maneras?
\opciones{12}{24}{120}{72}
\sol{B. Bloques: $2!$; dentro: $3! \cdot 2!$. $2 \cdot 6 \cdot 2 = 24$.}
\end{ejercicios}

\bitacora

% ------------------------------------------------------------------
\leccion{pu-co-3}{Combinaciones}{Juego de los grupos que se forman sin importar el orden.}

\subsection{La idea: el orden no importa}
Para elegir un \emph{comité} de 3 personas entre 5, el orden no importa: \{Ana, Beto, Caro\} es el mismo comité que
\{Caro, Ana, Beto\}. Si se cuentan los órdenes ($5 \cdot 4 \cdot 3 = 60$), cada comité aparece $3! = 6$ veces. Así que hay
$\tfrac{60}{6} = 10$ comités. Eso es una \clave{combinación}:
\[\binom{n}{k} = \frac{n(n - 1)\cdots(n - k + 1)}{k!} = \frac{n!}{k!\,(n - k)!}.\]

\subsubsection{¿Permutación o combinación?}
Pregúntate: si cambio el orden de los elegidos, ¿obtengo algo distinto? Cargos, podios, contraseñas: sí (permutación).
Comités, grupos, manos de cartas, partidas entre dos: no (combinación).

\subsubsection{Propiedades útiles}
$\binom{n}{k} = \binom{n}{n - k}$ (elegir los que entran es elegir los que salen). $\binom{n}{2} = \tfrac{n(n-1)}{2}$:
apretones de mano, partidas, diagonales.

\figura{ordenimporta}{Permutación y combinación: la misma gente, distinta pregunta.}

\begin{regla}
\begin{itemize}
  \item $\binom{n}{k} = \dfrac{P(n, k)}{k!}$. $\binom n0 = \binom nn = 1$, $\binom n1 = n$.
  \item $\binom nk = \binom n{n-k}$.
  \item Si hay dos grupos (hombres y mujeres), se elige de cada grupo y se multiplica.
\end{itemize}
\end{regla}

\subsection{Más ejemplos, en corto}
\begin{ejemplo}[Comité mixto]
Hay 6 hombres y 4 mujeres. ¿Cuántos comités de 3 hombres y 2 mujeres se pueden formar?
\solucion $\binom63 \cdot \binom42 = 20 \cdot 6 = 120$.
\end{ejemplo}

\begin{ejemplo}[Diagonales]
¿Cuántas diagonales tiene un octágono?
\solucion Segmentos entre 8 vértices: $\binom82 = 28$. Menos los 8 lados: 20.
\end{ejemplo}

\begin{ejemplo}[Al menos una mujer]
Con los mismos 10, ¿cuántos comités de 3 tienen al menos una mujer?
\solucion Total $\binom{10}{3} = 120$; sin mujeres $\binom63 = 20$. Con al menos una: 100.
\end{ejemplo}

\begin{trampa}
Usar permutaciones donde el orden no importa (el reto daría 90). Cada partida entre A y B se contaría dos veces: «A contra
B» y «B contra A».
\end{trampa}

\begin{reto}
En un torneo de ajedrez con 10 jugadores, cada uno juega exactamente una partida contra cada uno de los demás. ¿Cuántas
partidas se juegan?
\opciones{90}{45}{100}{20}
\end{reto}

\begin{retoresuelto}
Cada partida es una pareja de jugadores, sin orden: $\binom{10}{2} = \tfrac{10 \cdot 9}{2} = 45$. \textbf{Opción B.}
\end{retoresuelto}

\taller
\piso{Entrena}{elegir}
\begin{ejercicios}
\item $\binom{7}{3}$ y $\binom{7}{4}$. \sol{35 y 35.}
\item ¿Cuántos apretones de mano hay entre 12 personas si todos se saludan? \sol{66.}
\item ¿De cuántas maneras se eligen 2 sabores de helado de 8? \sol{28.}
\end{ejercicios}
\piso{Como en el examen}{opción múltiple}
\begin{ejercicios}
\item ¿Cuántos triángulos se forman con vértices en 7 puntos de una circunferencia?
\opciones{35}{21}{210}{42}
\sol{A. $\binom73$.}
\item De 5 libros de cuentos y 4 de poesía se eligen 3 de cuentos y 2 de poesía. ¿De cuántas maneras?
\opciones{60}{126}{20}{36}
\sol{A. $10 \cdot 6$.}
\end{ejercicios}

\bitacora

% ------------------------------------------------------------------
\leccion{pu-co-4}{Conteo por casos y complemento}{Juego de contar lo que sobra.}

\subsection{La idea: a veces es más fácil contar lo contrario}
Cuando una condición dice «al menos uno», «no todos» o «ni… ni…», suele ser más rápido contar \clave{lo que no cumple} y
restarlo del total. Y cuando una condición cambia según la situación, se separa en \clave{casos} que no se solapan y se
suman.

\subsubsection{Inclusión-exclusión}
Para contar los que cumplen $A$ o $B$: se suman y se resta lo contado dos veces, como en los diagramas de Venn:
$n(A \cup B) = n(A) + n(B) - n(A \cap B)$.

\begin{regla}
\begin{itemize}
  \item Complemento: casos buenos $=$ total $-$ casos malos.
  \item Casos: separa en situaciones excluyentes y suma.
  \item Múltiplos de $k$ entre 1 y $N$: la parte entera de $\tfrac Nk$.
\end{itemize}
\end{regla}

\subsection{Más ejemplos, en corto}
\begin{ejemplo}[Al menos un 7]
¿Cuántos números de tres cifras tienen al menos un 7?
\solucion Total: 900. Sin ningún 7: primera cifra 8 opciones (ni 0 ni 7), las otras 9 cada una: $8 \cdot 9 \cdot 9 =
648$. Con al menos un 7: $900 - 648 = 252$.
\end{ejemplo}

\begin{ejemplo}[Por casos]
¿Cuántos números de dos cifras tienen la suma de sus cifras igual a 10?
\solucion Primera cifra $a$ de 1 a 9 y $b = 10 - a$ debe ser cifra (0 a 9): $a$ de 1 a 9 da $b$ de 9 a 1. Son 9 números
(19, 28, …, 91).
\end{ejemplo}

\begin{trampa}
Al contar por casos, que dos casos se solapen y un mismo objeto se cuente dos veces.
\end{trampa}

\begin{reto}
¿Cuántos números del 1 al 100 no son divisibles ni entre 2 ni entre 3?
\opciones{33}{34}{17}{50}
\end{reto}

\begin{retoresuelto}
Divisibles entre 2: 50. Entre 3: 33. Entre los dos (entre 6): 16. Entre 2 o 3: $50 + 33 - 16 = 67$. Ni uno ni otro:
$100 - 67 = 33$. \textbf{Opción A.}
\end{retoresuelto}

\taller
\piso{Entrena}{complemento}
\begin{ejercicios}
\item ¿Cuántos números del 1 al 60 son múltiplos de 4 o de 6? \sol{$15 + 10 - 5 = 20$.}
\item Se lanza un dado 3 veces. ¿En cuántas secuencias sale al menos un 6? \sol{$216 - 125 = 91$.}
\end{ejercicios}
\piso{Como en el examen}{opción múltiple}
\begin{ejercicios}
\item ¿Cuántas palabras de 4 letras con las letras A, B, C (pueden repetirse) usan al menos una A?
\opciones{65}{81}{16}{64}
\sol{A. $3^4 - 2^4 = 81 - 16$.}
\item De 8 personas se eligen 3, pero Juan y Pedro no pueden estar juntos. ¿Cuántos grupos hay?
\opciones{56}{50}{36}{6}
\sol{B. Total 56; con los dos: 6 (el tercero entre los otros 6). $56 - 6$.}
\end{ejercicios}

\bitacora

% ------------------------------------------------------------------
\leccion{pu-co-5}{Contar figuras y caminos}{Juego de la cuadrícula de calles.}

\subsection{La idea: un camino es una palabra}
Cada camino corto es una secuencia de pasos: por ejemplo DDADA (derecha, derecha, arriba, derecha, arriba). Contar
caminos es contar palabras con 3 letras D y 2 letras A: $\tfrac{5!}{3!\,2!} = \binom52 = 10$.

\figura{caminos}{En cada esquina, el número de caminos que llegan a ella.}

El dibujo muestra otro método, útil cuando hay obstáculos: en cada esquina se escribe el número de caminos que llegan
a ella, que es la suma de los de la esquina de la izquierda y la de abajo (\clave{método de sumar esquinas}).

\subsubsection{Contar figuras}
Para contar triángulos, cuadrados o rectángulos en un dibujo, hay que \emph{organizar}: por tamaño, o por cuántas piezas
pequeñas forman cada figura. En una cuadrícula de $m \times n$, los rectángulos son $\binom{m+1}{2}\binom{n+1}{2}$ (se
eligen dos rectas verticales y dos horizontales).

\begin{regla}
\begin{itemize}
  \item Caminos en una cuadrícula de $a \times b$: $\binom{a + b}{a}$.
  \item Con obstáculos: sumar esquinas.
  \item Rectángulos en $m \times n$: $\binom{m + 1}{2}\binom{n + 1}{2}$. Cuadrados en $n \times n$: $1^2 + 2^2 + \cdots + n^2$.
\end{itemize}
\end{regla}

\subsection{Más ejemplos, en corto}
\begin{ejemplo}[Cuadrados en un tablero]
¿Cuántos cuadrados hay en una cuadrícula de $3 \times 3$?
\solucion De $1 \times 1$: 9. De $2 \times 2$: 4. De $3 \times 3$: 1. Total 14.
\end{ejemplo}

\begin{ejemplo}[Pasando por un punto]
En una cuadrícula de $4 \times 4$, ¿cuántos caminos cortos de la esquina inferior izquierda a la superior derecha pasan por
el centro $(2, 2)$?
\solucion Hasta el centro: $\binom42 = 6$. Del centro al final: 6. Total $36$.
\end{ejemplo}

\begin{trampa}
Contar figuras sin orden y olvidar las grandes (las formadas por varias piezas).
\end{trampa}

\begin{reto}
En una cuadrícula de calles, para ir de la esquina A a la esquina B hay que avanzar 3 cuadras hacia la derecha y 2 hacia
arriba, sin retroceder. ¿Cuántos caminos distintos hay?
\opciones{5}{6}{10}{12}
\end{reto}

\begin{retoresuelto}
Cada camino es una palabra con 3 D y 2 A: $\binom52 = 10$. La tabla de esquinas también termina en 10.
\textbf{Opción C.}
\end{retoresuelto}

\taller
\piso{Entrena}{contar}
\begin{ejercicios}
\item Caminos cortos en una cuadrícula de $4 \times 2$. \sol{$\binom62 = 15$.}
\item Rectángulos en una cuadrícula de $2 \times 3$. \sol{$\binom32 \binom42 = 3 \cdot 6 = 18$.}
\item ¿Cuántos triángulos hay en un triángulo dividido por dos rectas que salen de un vértice al lado opuesto? \sol{En el lado opuesto quedan 4 puntos: cada par forma un triángulo con el vértice: $\binom42 = 6$.}
\end{ejercicios}
\piso{Como en el examen}{opción múltiple}
\begin{ejercicios}
\item ¿Cuántos cuadrados hay en un tablero de $4 \times 4$?
\opciones{16}{25}{30}{20}
\sol{C. $16 + 9 + 4 + 1$.}
\item En la cuadrícula de la lección ($3 \times 2$), si la esquina $(1, 1)$ está cerrada, ¿cuántos caminos quedan?
\opciones{10}{6}{4}{8}
\sol{C. Caminos que pasan por $(1, 1)$: $\binom21 \cdot \binom31 = 2 \cdot 3 = 6$. Quedan $10 - 6 = 4$.}
\end{ejercicios}

\bitacora

\evaluacion
\begin{ejercicios}
\item ¿Cuántos números de tres cifras tienen todas sus cifras pares?
\opciones{100}{125}{64}{80}
\sol{A. Primera: 2, 4, 6, 8 (4); las otras 5 cada una: 100. (4.1)}
\item Con 4 camisas, 3 pantalones y 2 pares de zapatos, ¿cuántas pintas distintas hay?
\opciones{9}{24}{12}{18}
\sol{B. (4.1)}
\item ¿Cuántos anagramas tiene la palabra CASA?
\opciones{24}{12}{6}{4}
\sol{B. $\tfrac{4!}{2!}$. (4.2)}
\item ¿De cuántas maneras se ordenan 4 personas en una mesa redonda?
\opciones{24}{6}{12}{4}
\sol{B. $3!$. (4.2)}
\item ¿Cuántos grupos de 4 se forman con 9 personas?
\opciones{126}{3\,024}{36}{84}
\sol{A. $\binom94$. (4.3)}
\item Cada par de 15 ciudades está unido por una carretera. ¿Cuántas carreteras hay?
\opciones{105}{210}{225}{30}
\sol{A. $\binom{15}{2}$. (4.3)}
\item ¿Cuántos números de 2 cifras tienen al menos una cifra igual a 5?
\opciones{18}{19}{17}{20}
\sol{A. Total 90; sin 5: $8 \cdot 9 = 72$. (4.4)}
\item Caminos cortos en una cuadrícula de $5 \times 1$.
\opciones{5}{6}{10}{15}
\sol{B. $\binom61$. (4.5)}
\end{ejercicios}
